\section{Nonlinear magnetic data without pretending it is globally linear}
\label{sec:nonlinear}

Two kinds of parameter variation must be separated. Exogenous values such as
\begin{equation}
 R_s=R_s(T_s,f_e),\qquad R_{\mathrm{exc}}=R_{\mathrm{exc}}(T_r)
\end{equation}
may be updated from temperature and frequency estimates before the analytical
solve. At each update, the quadratic structure and all transformations remain
valid.

Saturation and cross-saturation are different. Writing
\(L_d=L_d(i_d,i_q,i_{\mathrm{exc}})\) inside a formula derived for a global
linear map mixes secant and differential quantities and generally destroys
the derivation. Start instead with a flux map
\begin{equation}
 \vect\psi=\gls{sym:phi}(\gls{sym:i}).
\end{equation}
At iterate \(\vect i_k\), a genuine tangent approximation is
\begin{equation}
 \vect\psi(\vect i)\approx
 \gls{sym:lk}\vect i+\gls{sym:psi0k},\qquad
 \matr L_k=\left.\frac{\partial\boldsymbol\phi}{\partial\vect i}
 \right|_{\vect i_k},\quad
 \vect\psi_{0,k}=\boldsymbol\phi(\vect i_k)-\matr L_k\vect i_k.
 \label{eq:local-affine-flux}
\end{equation}
The offset is essential: omitting it forces the tangent plane through the
origin and makes it fail even at the expansion point. Within one local affine
region, torque and voltage become affine-quadratic expressions and retain much
of the \gls{qcqp} structure. Recent work formulates saturated \gls{wrsm}
reference generation over \gls{pwa} finite-element flux regions in precisely
this spirit; as of this paper's date it is an arXiv preprint and is cited as
such \cite{parsonscherban2026}.

\begin{figure}[tb]
\centering
\begin{tikzpicture}[
 node distance=8mm and 9mm,
 block/.style={draw,rounded corners,align=center,minimum height=10mm,text width=.19\linewidth,fill=gray!6},
 flow/.style={-{Latex[length=2mm]},thick}]
 \node[block] (linear) {linear geometric\\initial solution};
 \node[block,right=of linear] (map) {read flux map\\$\matr L_k,\vect\psi_{0,k}$};
 \node[block,right=of map] (inner) {analytic local\\inner solve};
 \node[block,right=of inner] (check) {residual and\\region checks};
 \draw[flow] (linear)--(map);
 \draw[flow] (map)--(inner);
 \draw[flow] (inner)--(check);
 \draw[flow] (check.south) |- ++(0,-.8) -| node[pos=.25,below]{damp/update} (map.south);
\end{tikzpicture}
\caption{Saturation is handled by an outer flux-map update around the small
analytic geometric solve. The affine offset is retained, and convergence is
judged in the original nonlinear model.}
\label{fig:nonlinear-loop}
\end{figure}


A controller-oriented outer loop is:
\begin{enumerate}
 \item initialize with the global linear geometric solution;
 \item read \(\boldsymbol\phi(\vect i_k)\), \(\matr L_k\), loss parameters,
       and the correct affine offset from the map;
 \item solve the small local geometric allocation problem;
 \item damp the update,
       \(\vect i_{k+1}=(1-\gamma)\vect i_k+\gamma\vect i_{\mathrm{local}}\),
       if region switching or saturation makes the full step unreliable;
 \item stop when current change, torque residual, voltage residual, and flux
       model residual all meet tolerances.
\end{enumerate}
Two or three iterations are a design target, not a theorem. A contraction
check supports fixed-point convergence locally; an \gls{sqp} method offers
more systematic globalization but carries Hessian, line-search, and active-set
costs. The geometric inner solution is valuable in both cases as a warm start,
preconditioner, and diagnostic. If the local model crosses a PWA boundary,
candidate neighboring regions must be checked rather than assuming one affine
piece remains valid.

\paragraph{What survives nonlinearity?}
Global ellipsoids and exact homogeneity do not. Local loss metrics, directional
voltage gains, active-flux combinations, residual checks, and the separation
between an inner allocation solve and an outer map update do. That is the
structure worth preserving.
